\naslov{Lagrangeova metoda}

\begin{definicija}
  Realna funkcija $\phi(x, y, z)$ se imenuje \pojem{prvi integral} $\zvodvi{1}$
  polja $V = (a,b,c)$, če je konstantna vzdolž vsake krivulje $\gamma(t)$, ki
  reši $\dot{\gamma} = V(\gamma)$.
\end{definicija}

\vprasanje{Kaj je prvi integral?}

Velja
\[
  \partial_t \phi(\gamma(t)) = \sk{\grad \phi(\gamma(t)), \dot{\gamma(t)}},
\]
torej je $\phi$ prvi integral za $V$ natanko tedaj, ko je za vsako tokovnico
$\gamma$ za $V$ ta skalarni produkt enak $0$.

\begin{trditev}
  Če $u$ reši kvazilinearno enačbo $a u_x + b u_y = c$, je $\psi(x,y,z) = z -
  u(x,y)$ prvi integral za prirejeno polje $V = (a,b,c)$.
\end{trditev}

\begin{proof}
  Izračunamo $\sk{\grad \psi, V} = -au_x - b u_y + c = 0$.
\end{proof}

\begin{trditev}
  Vsaka nivojna ploskev $\Sigma = \{\phi = C\}$ prvega integrala $\phi$ je unija
  karakteristik polja $V$.
\end{trditev}

\begin{proof}
  Vzemimo točko $p_0 \in \Sigma$ in rešimo Cauchyjevo nalogo $\dot{\gamma} =
  V(\gamma), \gamma(0) = p_0$.
  Tako dobimo tokovnico, ki vsebuje $p_0$.
  Obratno, tokovnica, ki seka $\Sigma$ v točki $t_0$, cela leži v $\Sigma$, saj
  je $\phi$ prvi integral.
\end{proof}

Integralne ploskve za kvazilinearno enačbo so torej nivojnice prvih integralov
polja $V = (a,b,c)$.
Nivojnica prvega integrala pa je integralna ploskev samo, če na njej lahko
izrazimo $z$ kot funkcijo spremenljivk $x, y$.

\begin{trditev}
  Naj bo $\phi \in \zvodvi{1}$ prvi integral za $V = (a,b,c)$ v okolici točke
  $p_0 \in \R^3$.
  Označimo $d = \phi(p_0)$.
  Če je $\partial_z \phi(p_0) \ne 0$, je nivojnica $S = \{\phi = d\}$ v okolici
  integralna ploskev za kvazilinearno enačbo.
\end{trditev}

\begin{proof}
  Ker je odvod neničeln, je $S$ lokalno graf neke $\zvodvi{1}$ funkcije $u$.
  Dokazati želimo, da $u$ reši kvazilinearno enačbo.
  Lokalno velja $\phi(x,y,u(x,y)) = d$; če to odvajamo po $x$ in $y$, dobimo
  \begin{gather*}
	\sk{\grad \phi, (1, 0, u_x)} = 0, \\
	\sk{\grad \phi, (0, 1, u_y)} = 0.
  \end{gather*}
  Torej je $\grad \phi$ vzporeden vektorskemu produktu $(1,0,u_x) \times
  (0,1,u_y) = (-u_x, -u_y, 1)$.
  Ker pa je $\phi$ prvi integral polja $V = (a,b,c)$, je nanj pravokoten, torej
  sta vektorja $(a,b,c)$ in $(-u_x, -u_y, 1)$ pravokotna, kar je ravno
  kvazilinearna enačba.
\end{proof}

\vprasanje{Dokaži, da so nivojnice prvega integrala polja $(a,b,c)$ integralne
  ploskve za kvazilinearno enačbo $a u_x + b u_y = c$.}

\begin{definicija}
  Če je $\phi$ prvi integral za $V$ in $d \in \R$, se družina nivojnic $\{\phi =
  d\}$ imenuje \pojem{splošna rešitev} za kvazilinearno enačbo $a u_x + b u_y =
  c$.
\end{definicija}

\begin{izrek}
  Naj bosta $\phi_1$ in $\phi_2$ neodvisna prva integrala polja $(a,b,c)$ v
  okolici točke $p \in \R^3$.
  Splošna rešitev $a u_x + b u_y = c$ je v okolici točke $p$ podana z
  $\{F(\phi_1, \phi_2) = 0\}$, kjer je $F = F(x,y)$ poljubna funkcija razreda
  $\zvodvi{1}$.
\end{izrek}

\begin{proof}
  Preverimo, da je $F \circ (\phi_1, \phi_2) = 0$ res posplošena rešitev.
  Če je $\gamma$ tokovnica in je $\phi_1 = d_1$, $\phi_2 = d_2$ na $\gamma$, je
  $F(\phi_1, \phi_2) \circ \gamma$ konstanta.

  V drugo smer preverimo, da je rešitev res take oblike.
  Vzemimo $p_0 \in \R^3$, ki leži v definicijskem območju $V$.
  Vemo, da je $\sk{\grad \phi(p_0), V(p_0)} = 0$ za vsak prvi integral $\phi$.
  Naj bo $\phi_3$ še en prvi integral.
  Tedaj je $V(p_0) \in \jedro D(\phi_1, \phi_2, \phi_3)$, ker je pravokoten na
  vse gradiente.
  Hkrati je $V \ne 0$ povsod, ker bi sicer nekje veljalo $\dot{\gamma} = 0$.
  Torej je Jacobijeva matrika izrojena in so funkcije $\phi_1$, $\phi_2$ in
  $\phi_3$ odvisne, zato obstaja neničelna funkcija $G = G(x,y,z)$, da lokalno
  velja $G(\phi_1, \phi_2, \phi_3) = 0$.
  Nivojnica $\{\phi_3 = d\}$ je lokalno vsebovana v $\{G(\phi_1, \phi_2,
  \phi_3) = 0\}$, kar je natanko $\{F(\phi_1, \phi_2) = 0\}$, če za $F$ vzamemo
  $F(x,y) = G(x,y,d)$.
\end{proof}

\vprasanje{Utemelji pravilnost metode prvih integralov.}

\begin{izrek}
  Naj bo $V$ vektorsko polje razreda $\zvodvi{1}$ v okolici točke $p \in \R^3$
  in naj bo $V(p) \ne 0$.
  Tedaj obstaja $\zvodvi{1}$ vektorsko polje $\psi$, definirano na okolici točke
  $p$, da je $\sk{\grad \psi_j, V}_{\R^3} = \delta_{1j}$ v okolici $p$.
\end{izrek}

\begin{proof}
  Privzemimo prvo $p = 0$ in $V(p) = (1,0,0)$.
  Naj bo $\phi_t$ tok polja $V$, tj.~družina $\zvodvi{1}$ vektorskih polj, za
  katera je $\phi_t(x) = V(\phi_t(x))$ in $\phi_0(x) = x$.
  V okolici točke $0$ definiramo še $g(x_1, x_2, x_3) = \phi_{x_1}(0, x_2,
  x_3)$.
  Dobimo
  \[
	\partial_{x_1} g(x)
	= \lim_{h \to 0} \frac{\phi_{x_1 + h}(0, x_2, x_3) - \phi_{x_1}(0, x_2,
	  x_3)}{h}
	= \partial_{x_1} \phi((0,x_2, x_3), x_1)
	= V(g(x))
  \]
  za $\phi(x,t) = \phi_t(x)$.
  Velja še $g(0,x_2, x_3) = (0, x_2, x_3)$, zato je $\partial_{x_j} g = e_j$ na
  ravnini $x_1 = 0$.
  Iz $V(g(0)) = v(0) = (1,0,0)$ sledi $Dg(0) = I_3$, zato je $g$ difeomorfizem
  na neki okolici točke $0$.
  Označimo $\psi = g^{-1}$.
  Če pišemo $x = (x_1, x_2, x_3)$ in $\psi = (\psi_1, \psi_2, \psi_3)$, tedaj iz
  zveze $\psi \circ g = \id$ sledi $\psi_j(g(x)) = x_j$.
  Odvajamo po $x_1$ in dobimo
  \[
	\delta_{1j}
	= \sum_{k=1}^3 \partial_{y_k} \psi_j(g(x)) \partial_{x_1} g_k(x)
	= \sum_{k=1}^3 \partial_{y_k} \psi_j(g(x)) V_k(g(x))
	= \sk{\grad \psi_j(g(x)), V(g(x))}.
  \]
  Ker je $g$ difeomorfizem, velja $\sk{\grad \psi_j, V} = 0$ na neki okolici
  $0$.

  Sedaj obravnavajmo splošni primer.
  Če sta $p, V(p) \in \R^3$ poljubna, vzamemo ortogonalno $A \in \R^{3 \times
	3}$, s katero se $\normiran{V(p)}$ slika v $(1,0,0)$.
  Označimo $\lambda = \norm{V(p)}^{-1}$ in $w(x) = \lambda A(V(x + p))$.
  Ta $w$ slika iz okolice točke $0$, velja $w(0) = (1,0,0)$.
  Torej obstaja vektorsko polje $\eta$, definirano na okolici $0$, da je
  $\sk{\grad \eta_j, w} = \delta_{1j}$.
  Sedaj definiramo še $\psi$ tako, da velja $\eta(x) = \lambda^{-1} A \psi(x +
  p)$.
  Izračunamo lahko, da velja
  \[
	\sk{\grad \eta_j(x), w(x)} = \sk{\grad \psi_j(x + p), V(x+p)}.
  \]
\end{proof}

\vprasanje{Dokaži, da za vsako vektorsko polje $V$ in točko $p$, za katero $V(p)
  \ne 0$, obstaja polje $\psi$, da je v okolici $p$ skalarni produkt $\sk{\grad
	\psi_j, V} = \delta_{1j}$.}